% !PRINTED 208
\weberpage{208}
%<<CONT
durch eine andere ersetzt wird, indem man in (5) $\nu'$ durch ein anderes Zeichen, $m$, ersetzt und dann $nu = 2\nu'\omega$ setzt. Beschränken wir uns auf den Fall $(g_1, g_2) = (0, 1)$ und multipliziert (5) noch mit $\vartheta_{00} : \vartheta_{10}$, so folgt: % [sic] Source grammar prints "multipliziert"; "multiplizieren wir" is intended.

\weq{12}{
\sum^{\nu} e^{\frac{8m\nu\pi i}{n}} x_{\nu,\nu'} = \\
(-1)^{\frac{n-1}{2}} n\, e^{\frac{-4\pi i m\nu'\omega}{n}} \frac{\vartheta_{00}\,\vartheta_{01}\,\vartheta'_{11}(0, n\omega)\,\vartheta_{11}[2(\nu'-m)\omega, n\omega]}{\vartheta_{10}\,\vartheta'_{11}\,\vartheta_{01}(2m\omega, n\omega)\,\vartheta_{01}(2\nu'\omega, n\omega)},
}
% Source verified: the summation sign carries only the upper index \nu.

wo die linke Seite \gesperrt{dann}, aber auch \gesperrt{nur} dann verschwindet, wenn $m \equiv \nu' \pmod{n}$ ist.

\medskip

\begin{center}
\textbf{\S\,63. Die \gesperrt{Galoissche Gruppe} der Teilungsgleichung.}
\end{center}

Im dreizehnten Abschnitt des ersten Bandes ist gezeigt, wie die algebraische Natur einer Gleichung ihren einfachsten Ausdruck in der \gesperrt{Galoisschen Gruppe} der Gleichung findet.

Es sei hier kurz an die Definition der \gesperrt{Galoisschen Gruppe} erinnert.

Nachdem der Rationalitätsbereich $\Omega$ festgesetzt war, ist die \gesperrt{Galoissche Gruppe} einer Gleichung $F(x) = 0$, von der nur vorausgesetzt wird, daß sie keine gleichen Wurzeln habe, definiert als die Gruppe $G$ der Permutationen unter den Wurzeln dieser Gleichung, der die doppelte Eigenschaft zukommt:

a) Jede rationale Gleichung in $\Omega$, die zwischen den Wurzeln von $F(x)$ besteht, bleibt richtig, wenn die Wurzeln irgend einer Permutation dieser Gruppe unterworfen werden.

b) Jede rationale Funktion in $\Omega$ von den Wurzeln von $F(x)$, die sämtliche Permutationen der Gruppe gestattet, ist in $\Omega$ enthalten.

Indem wir nun die Gruppe $G$ der Teilungsgleichung zu bestimmen suchen, setzen wir als Rationalitätsbereich zunächst den \gesperrt{Inbegriff aller rationalen Zahlen und rationalen Funktionen von $\varkappa^2$ voraus}.

Nach dem Additions- und Multiplikationstheorem (vgl. \S\, 57, 61) ist, wenn $f$ und $\varphi_m$ rationale Funktionen bedeuten, die von $\mu, \mu' \nu, \nu'$ unabhängig sind, % [sic] The source omits the comma between \mu' and \nu.

\weq{1}{x_{\mu+\nu,\,\mu'+\nu'} = f(x_{\mu,\mu'}, x_{\nu,\nu'}),}

\weq{2}{x_{m\mu,\,m\mu'} = \varphi_m(x_{\mu,\mu'}).}
